\selectlanguage{ngerman}%
\frenchspacing%
\begin{ledgroupsized}[r]{120mm}%
\footnotesize%
\pstart%
\noindent%
\textbf{Überlieferung:}%
\pagelabel{Gerland_1906_007-009_uebrlf}%
\pend%
\end{ledgroupsized}%
\begin{ledgroupsized}[r]{114mm}%
\footnotesize%
\pstart%
\parindent -6mm%
\makebox[6mm][l]{\textit{E}}%
Druck nach unbekannter Vorlage:
\textsc{Gerland} 1906, S.\;7--9 (Nr.\;5).\cite{00197}
Die verschollene handschriftliche Vorlage wies laut \textit{E} (S.\;7) das Format eines Oktav-Blattes auf.
Der oft fehlerhafte Text von \textit{E} %
wird hier wiedergegeben und ggf. berichtigt oder kommentiert;
die Rechtschreibung wird an die Editionsregeln angepasst.
\pend%
\end{ledgroupsized}%
\vspace{5mm}%
\begin{ledgroup}%
\footnotesize%
\pstart%
\noindent%
\textbf{Datierungsgründe:}
Die Datierung des Konzepts N.\;\ref{56215} beruht auf inhaltlichen Betrachtungen.
Siehe die ausführliche Begründung auf S.\;\pageref{56215_dat}.%
\pend%
\end{ledgroup}%
\selectlanguage{latin}%
\frenchspacing%
\vspace{6mm}%
\normalsize%
\pstart%
\noindent%
\lbrack S.\;7\rbrack\
\hfill%
Aero-nautica%
\indexS{aeronautica}
\hfill%
\phantom{\lbrack S.\;7\rbrack}%
\pend%
\vspace{0.5em}
\pstart%
\noindent%
\edtext{%
\edlabel{Gerland_1906_007-009_hestz-1}%
Aeris pes cubicus pondere aequat unciam unam cum dimidia.%
\edlabel{Gerland_1906_007-009_hestz-2}%
\indexS{uncia}%
\indexS{pes cubicus aeris}%
\indexS{pondus aeris communis}%
\indexS{aer communis}%
}{%
\lemma{Aeris \lbrack...\rbrack\ dimidia}%
\Cfootnote{%
Vgl. P.\:\textsc{Lohmeier}, \textit{Exercitatio physica}, §\;8 (Rinteln 1676, S.\;7f.).\cite{01422}
Siehe zudem F.\,\textsc{Lana\:de\:Terzi},
\textit{Prodromo}, I Teil, cap.\;6 (Brescia 1670, S.\;53).\cite{00069}%
}}
\edtext{%
\edlabel{Gerland_1906_007-009_xtd-1}%
Laminae e cupro satis adhuc consistentes%
\indexS{lamina cuprea}%
\indexS{lamina consistens}
pedem unum longae et latae,%
\indexS{pes}%
\indexS{pes quadratus laminae cupreae}
quae tres uncias suo pondere non excedant;%
\indexS{pondus laminae cupreae}%
\indexS{uncia}%
\edlabel{Gerland_1906_007-009_xtd-2}%
}{%
\lemma{Laminae \lbrack...\rbrack\ excedant}%
\Cfootnote{Vgl. \textsc{Lohmeier}, \textit{Exercitatio physica}, §\;21 (S.\;14) u. §\;30 (S.\;18).\cite{01422}
Siehe zudem \textsc{Lana\:de\:Terzi},
\textit{Prodromo}, I Teil, cap.\;6 (S.\;56).\cite{00069}%
}}
\edtext{%
\edlabel{Gerland_1906_007-009_crj-1}%
quarum 12 constituunt libram.%
\indexS{libra}%
\edlabel{Gerland_1906_007-009_crj-2}%
}{%
\lemma{quarum \lbrack...\rbrack\ libram}%
\Cfootnote{Vgl. \textsc{Lohmeier}, \textit{Exercitatio physica}, §\;8 (S.\;8) u. §\;21 (S.\;14).\cite{01422}
Siehe zudem \textsc{Lana\:de\:Terzi},
\textit{Prodromo}, I Teil, cap.\;6 (S.\;53).\cite{00069}%
}}
Sphaera ex his laminis,%
\indexS{sphaera e lamina cuprea}
quae pedem haberet in diametro,%
\indexS{diameter sphaerae}%
\indexS{pes}
esset (secluso aere) ponderis $9\dfrac{3}{7}$ unciarum.%
\indexS{pondus sphaerae cupreae}%
\rule[0.0mm]{0.0mm}{5.6mm}
Nam diam. sphaerae est \textit{a},
nempe 1 pes,
circumf. $\dfrac{22}{7}a$,
superficies sphaerae est $\dfrac{22}{7}a^2$,%
\indexS{superficies sphaerae}%
\rule[0.0mm]{0.0mm}{5.6mm}
sit $a^2$
\edtext{%
pes%
\indexS{pes}%
}{%
\lemma{pes}%
\Cfootnote{%
Ein Quadratfuß
(vgl. \refpassage{Gerland_1906_007-009_xtd-1}{Gerland_1906_007-009_xtd-2}.)%
}}
\edtext{%
un\lbrack us\rbrack%
}{%
\lemma{una}%
\Bfootnote{%
\textit{E\;ändert Hrsg.}\cite{00197}%
}%
\lemma{un\lbrack us\rbrack}%
\Cfootnote{%
In \textit{E} ist die möglicherweise von einer Verlesung herrührende Unstimmigkeit \textit{una vel 3 unc.} mit \lbrack!\rbrack\ markiert.%
}%
}
vel 3 unc., %
erit superficies sphaerae%
\indexS{superficies sphaerae}
$\dfrac{66}{7}$ unc.%
\rule[0.0mm]{0.0mm}{5.6mm}
seu $9\dfrac{3}{7}$ unc.%
\indexS{uncia}
\pend%
\pstart%
Aeris inclusi pondus esset $\dfrac{11}{14}$ unc.%
\indexS{aer inclusus sphaerae cupreae}%
\indexS{pondus aeris inclusi}%
\indexS{uncia}
Nam
\edtext{%
area externa%
}{%
\lemma{area externa}%
\Cfootnote{%
Gemeint ist das Volumen der Kugel.
In \textit{E} ist die Textstelle mit \lbrack?\rbrack\ markiert.
Die Unstimmigkeit rührt möglicherweise von einer Verlesung her.
}}
sphaerae fit
multiplicata superficie eius per sextam partem diametri,%
\indexS{diameter sphaerae}
seu $\dfrac{22}{7 \cdot 6}a^3$,
est autem $a^3$ pondus $1\dfrac{1}{2}$ unciae,
fiet $\dfrac{22}{7 \cdot 6}a^3$ aequ. $\dfrac{11}{14}$ unc.%
\indexS{uncia}%
\rule[0.0mm]{0.0mm}{5.6mm}
\lbrack S.\;8\rbrack\
\pend%
\newpage%
\pstart%
Multiplicemus \textit{a} per \textit{b}
fiet $\dfrac{22}{42} a^3 b^3$ aequ. $\dfrac{11}{14}b^3$
et
\edtext{%
$\dfrac{22}{7}a^2[b^2]$%
}{%
\lemma{$b^2$}%
\Bfootnote{%
\textit{erg. Hrsg.}%
}}
aequ. $\dfrac{66}{7}b^2$,
debent esse $\dfrac{66}{7}b^2$ aequ. $\dfrac{11}{14}b^3$
seu fiet: 12 aequ. \textit{b},
itaque si ponatur sphaera pedum 12%
\indexS{pes}
aequale erit pondus sphaerae cupreae%
\indexS{pondus sphaerae cupreae}%
\indexS{sphaera cuprea}%
\rule[0.0mm]{0.0mm}{4.8mm}
et aeris ei inclusi.%
\indexS{pondus aeris inclusi}%
\indexS{aer inclusus sphaerae cupreae}
\pend%
\pstart%
Generaliter sphaerae cupreae pondus est%
\indexS{pondus sphaerae cupreae}%
\indexS{sphaera cuprea}
\edtext{%
\lbrack$\dfrac{66}{7}b^2$\rbrack%
}{%
\lemma{$\dfrac{11}{14}b^3$}%
\Bfootnote{%
\textit{E\;ändert Hrsg.}\cite{00197}%
}}
unciarum%
\indexS{uncia}%
\rule[0.0mm]{0.0mm}{5.6mm}
posito \textit{b} esse numerum pedum diametri,%
\indexS{diameter sphaerae}
et aeris inclusi pondus est%
\indexS{pondus aeris inclusi}%
\indexS{aer inclusus sphaerae cupreae}
\edtext{%
\lbrack$\dfrac{11}{14}b^3$\rbrack%
}{%
\lemma{$\dfrac{66}{7}b^2$}%
\Bfootnote{%
\textit{E\;ändert Hrsg.}\cite{00197}%
}}
unciarum%
\indexS{uncia}%
\rule[0.0mm]{0.0mm}{4.8mm}
posito \textit{b} numero eodem.%
\edlabel{Gerland_1906_007-009_scheitel-1}
\pend%
\pstart%
\edlabel{Gerland_1906_007-009_scheitel-2}%
Sed mea sententia%
\indexS{sententia}%
\rule[0.0mm]{0.0mm}{4.6mm}
laminae istae cupreae%
\indexS{lamina cuprea}
sunt nimis tenues.
Itaque calculum instituamus%
\indexS{calculus generalior}%
\rule[0.0mm]{0.0mm}{5.2mm}
generalius.
Numerus pedum diametri sphaerae%
\indexS{diameter sphaerae}%
\indexS{pes}
sit \textit{b}.
\edlabel{Gerland_1906_007-009_itu-1}%
\textit{A} sit 4%
\edlabel{Gerland_1906_007-009_itu-2}
\edtext{%
pe\lbrack de\rbrack s%
}{%
\lemma{pes}%
\Bfootnote{%
\textit{E\;ändert Hrsg.}\cite{00197}%
}}
et area superficiei sphaerae%
\indexS{superficies sphaerae}%
\indexS{sphaera cuprea}
$\dfrac{22}{7}b^2 4^2$
et area solidi%
\indexS{area solidi}%
\indexS{solidum sphaerae}
$\dfrac{11}{21}b^3 A^3$,%
\rule[0.0mm]{0.0mm}{5.6mm}
seu unius pedis quadrati cuprei pondus%
\indexS{pondus pedis quadrati cuprei}%
\indexS{pes quadratus}
esto~\textit{c},
similiter
\edtext{%
$A^3$ seu unius pedis cubici pondus%
\indexS{pes cubicus}%
\indexS{pondus pedis cubicis aeris}%
\indexS{pondus aeris inclusi}%
}{%
\lemma{$A^3$ \lbrack...\rbrack\ pondus}%
\Cfootnote{%
Die Angabe widerspricht der Setzung $A = 4$ Fuß
(vgl. \refpassage{Gerland_1906_007-009_itu-1}{Gerland_1906_007-009_itu-2}).%
}}
esto aequale $1\dfrac{1}{2}$ unc.%
\indexS{uncia}%
\rule[0.0mm]{0.0mm}{5.2mm}
\par%
\begin{edtabularl}%
&
fiet 
&
\hspace*{1.0mm}%
pondus cupreae sphaerae%
\indexS{pondus sphaerae cupreae}%
\indexS{sphaera cuprea}%
\rule[0.0mm]{0.0mm}{5.2mm} &
\hspace*{12.5mm}%
$\dfrac{22}{7}b^2 c$ & \\
& 
&
\hspace*{1.0mm}%
pondus autem aeris inclusi%
\indexS{pondus aeris inclusi}%
\indexS{aer inclusus sphaerae cupreae}%
\rule[0.0mm]{0.0mm}{5.6mm}%
&
\hspace*{12.5mm}%
$\dfrac{11}{21}b^3 \dfrac{3}{2}$ unc. & \\
& & &
\rule[0.0mm]{0.0mm}{5.6mm}%
\hspace*{13.0mm}%
seu $\dfrac{11}{14}b^3$ unc. &
\end{edtabularl}%
\pend%
\pstart%
\noindent%
Rem ergo breviter ita ponamus%
\rule[0.0mm]{0.0mm}{4.8mm}%
\lbrack:\rbrack\
\par%
\begin{edtabularl}%
&
Sit%
\rule[0.0mm]{0.0mm}{5.8mm} &
pondus pedis quadrati laminae cupreae%
\indexS{pondus laminae cupreae}%
\indexS{pes quadratus laminae cupreae} &
\textit{c} & \\
& &
varians pro crassitie laminae%
\indexS{crassities laminae cupreae}%
\indexS{lamina cuprea} & &
\\
& &
\rule[0.0mm]{0.0mm}{5.2mm}%
pondus%
\indexS{pondus pedis cubicis aeris}%
\indexS{pes cubicus aeris}%
\indexS{aer inclusus sphaerae cupreae}
\edtext{%
\lbrack pedis\rbrack%
}{%
\lemma{pedis}%
\Bfootnote{%
\textit{erg. Hrsg.}%
}}
cubici aeris &
$\dfrac{3}{2}$ unc. & \\
& &
\rule[0.0mm]{0.0mm}{4.2mm}%
numerus pedum diametri sphaerae%
\indexS{diameter sphaerae} &
\textit{b} & \\
&
Erit\rule[0.0mm]{0.0mm}{6.0mm} &
pondus sphaerae cupreae%
\indexS{pondus sphaerae cupreae}%
\indexS{sphaera cuprea} &
$\dfrac{22}{7}b^2 c$ & \\
\edatleft{\{}{4.0\baselineskip}
& &
pondus aeris inclusi%
\indexS{pondus aeris inclusi}%
\indexS{aer inclusus sphaerae cupreae}%
\rule[0.0mm]{0.0mm}{5.2mm} &
$\dfrac{11}{4}b^3$ unc. &
\edatright{\}}{4.0\baselineskip}%
\end{edtabularl}%
\pend%
\vspace{0.5em}%
\newpage%
\pstart%
\noindent%
Quare %
si laminam pedis quadrati cuprei%
\indexS{lamina cuprea}%
\indexS{pes quadratus laminae cupreae}
ponamus esse
\edlabel{Gerland_1906_007-009_crj-3}%
\edtext{%
duarum librarum%
\indexS{libra}
seu 32 unciarum:%
\indexS{uncia}%
}{%
\lemma{duarum \lbrack...\rbrack\ unciarum}%
\Cfootnote{%
Oben (\refpassage{Gerland_1906_007-009_crj-1}{Gerland_1906_007-009_crj-2}) wurde angegeben, dass ein Pfund zwölf Unzen gleiche.
Laut \textit{E} %
(S.\;9)\cite{00197} entspricht diese Einteilung dem "Apothekergewicht", während die Einteilung des Pfundes in sechzehn Unzen auf das "altfranzösische Gewicht" zurückgehe.
In der Frühneuzeit waren beide Systeme im deutschsprachigen Raum in Gebrauch.
Die Einteilung des Pfundes in sechzehn Unzen wird im Folgenden als Grundlage für die Berechnung verwendet.%
}}%
\edlabel{Gerland_1906_007-009_crj-4}
quaeratur magnitudo diametri talis,%
\indexS{magnitudo}%
\indexS{diameter sphaerae}
ut tantum ponderet sphaera cuprea%
\indexS{sphaera cuprea}%
\indexS{pondus sphaerae cupreae}
quantum aer inclusus,%
\indexS{aer inclusus sphaerae cupreae}%
\indexS{pondus aeris inclusi}
fiet:
\pend%
\pstart%
\noindent%
\centering%
$\dfrac{22}{7}b^2 \, 32 \ \text{unc.} \ \sqcap\, \ \dfrac{11}{14}b^3 \ \text{unc.}$
\pend%
\pstart%
\noindent%
seu fiet
$32 A \ \text{aequ.} \ b$%
\rule[0.0mm]{0.0mm}{4.6mm}
sive deberet esse sphaera 128 pedum.%
\indexS{pes}
\pend%
\pstart%
Generaliter
si quaeratur longitudo diametri sphaerae,%
\indexS{diameter sphaerae}%
\indexS{sphaera cuprea}
ut fiat aequatio,%
\indexS{aequatio}
habebitur $\dfrac{22}{7} b^2 c$ aequ. $\dfrac{11}{14}b^3$ unc.
seu 4\textit{c} unc. aequ. \textit{b},%
\indexS{uncia}
sive
\edtext{%
\edlabel{Gerland_1906_007-009_error-1}%
numerus \textit{b} quaesitus%
\indexS{numerus quaesitus}
erit ad~4 %
\edtext{%
\lbrack u\rbrack t%
}{%
\lemma{Ut}%
\Bfootnote{%
\textit{E\;ändert Hrsg.}\cite{00197}%
}}
pondus \textit{c} laminae pedis quadrati materiae,%
\indexS{pondus laminae cupreae}%
\indexS{pes quadratus laminae cupreae}%
\indexS{materia laminae cupreae}%
\indexS{lamina cuprea}
ex qua globum conficere volumus,%
\indexS{globus cupreus}
est ad unciam%
\indexS{uncia}%
\lbrack;\rbrack%
\edlabel{Gerland_1906_007-009_error-2}%
}{%
\lemma{numerus \lbrack...\rbrack\ unciam}%
\Cfootnote{%
Vielmehr gilt die Proportion:
$b : 4 = c\,\text{unc.} : 1.$
Vgl. zur Bestätigung den abschließenden Vermerk \refpassage{Gerland_1906_007-009_schlss-1}{Gerland_1906_007-009_schlss-2}.%
}}
itaque si lamina%
\indexS{lamina cuprea}
esset unius librae,%
\indexS{libra}
erit sphaera necessaria%
\indexS{sphaera cuprea}
minimum diametri 64 pedum%
\indexS{diameter sphaerae}%
\indexS{pes}%
\lbrack;\rbrack\
itaque rem ad usum transferri%
\indexS{res translata ad usum}%
\indexS{usus}
posse non puto.
Nam nec tantas sphaeras,%
\indexS{sphaera cuprea}
etiam per partes si quis parabit,
et si paratae essent,
quis credat laminas,%
\indexS{lamina cuprea}
si tanto ponderi comparentur tam tenues,
resistere posse ponderi tam immensi,%
\indexS{pondus immensum}
quantum est alias inviis ventis.%
\indexS{ventus invius}
Nec aequabilitas sphaerae%
\indexS{aequabilitas sphaerae cupreae}%
\indexS{sphaera cuprea}
\edtext{%
pr\lbrack aestari\rbrack%
}{%
\lemma{privare}%
\Bfootnote{%
\textit{E\;ändert Hrsg.}\cite{00197}%
}}
potest%
\lbrack:\rbrack\
ea enim tanta %
non potest
quia exiguae inaequalitates%
\indexS{inaequalitas exigua}
\edtext{%
supers\lbrack u\rbrack nt,%
}{%
\lemma{supersint}%
\Bfootnote{%
\textit{E\;ändert Hrsg.}\cite{00197}%
}}
quae horribili isti ponderi,%
\indexS{pondus horribile}
cui nihil in natura%
\indexS{natura}
simile visum est,%
\indexS{nihil simile in natura}
necessario cedent%
\lbrack;\rbrack\
itaque rem puto esse supra
\lbrack S.\;9\rbrack\
vires humanas.%
\indexS{res putanda supra vires humanas}%
\indexS{vis humana}
\pend%
\pstart%
\edlabel{Gerland_1906_007-009_schlss-1}%
\edtext{}{%
\xxref{Gerland_1906_007-009_schlss-1}{Gerland_1906_007-009_schlss-2}%
\lemma{Numerus \lbrack...\rbrack\ manentis}%
\Cfootnote{%
Wohl ein nachträglicher Vermerk, der die ungenaue Formulierung im Haupttext,
\refpassage{Gerland_1906_007-009_error-1}{Gerland_1906_007-009_error-2},
und in der Randbemerkung,
\refpassage{Gerland_1906_007-009_rndbmrk-1}{Gerland_1906_007-009_rndbmrk-2},
berichtigt.
Siehe die entsprechenden Erläuterungen.%
}}%
Numerus unciarum ponderis laminae pedis quadrati materiae%
\indexS{uncia}%
\indexS{pondus laminae cupreae}%
\indexS{pes quadratus laminae cupreae}%
\indexS{materia laminae cupreae}%
\indexS{lamina cuprea}
unde fit sphaera%
\indexS{sphaera e lamina cuprea}
quadruplicatus%
\lbrack,\rbrack\
est
\edtext{%
numer\lbrack us\rbrack%
}{%
\lemma{numero}%
\Bfootnote{%
\textit{E\;ändert Hrsg.}\cite{00197}%
}}
pedum diametris sphaerae in suspenso in%
\indexS{pes}%
\indexS{diameter sphaerae}%
\indexS{sphaera suspensa in aere}%
\indexS{aer communis}
\edtext{%
aer\lbrack e\rbrack%
}{%
\lemma{aeri}%
\Bfootnote{%
\textit{E\;ändert Hrsg.}\cite{00197}%
}}
manentis.%
\edlabel{Gerland_1906_007-009_schlss-2}%
\pend%
\vspace{1.5em}%
\pstart%
\noindent%
\lbrack%
\textit{In der verschollenen handschriftlichen Vorlage war die folgende Randbemerkung laut E (S.\;9) \glqq in die Überschrift hineingeschrieben\grqq:}%
\rbrack%
\pend%
\vspace{0.5em}%
\pstart%
\noindent%
\edlabel{Gerland_1906_007-009_rndbmrk-1}%
\edtext{%
(Conclusi
numerum pedum diametri sphaerae esse ad numerum 4,%
\indexS{diameter sphaerae}
ut numerus unciarum ponderis laminae pedis quadrati
unde fieri debet sphaera,%
\indexS{sphaera e lamina cuprea}
est ad unciam%
\lbrack;\rbrack%
\edlabel{Gerland_1906_007-009_rndbmrk-2}%
}{%
\lemma{(Conclusi \lbrack...\rbrack\ unciam}%
\Cfootnote{%
Wieder die ungenaue Formulierung wie in der Passage
\refpassage{Gerland_1906_007-009_error-1}{Gerland_1906_007-009_error-2}.
Die Fortsetzung der Randbemerkung berichtigt die Aussage.%
}}
\edlabel{Gerland_1906_007-009_rndbmrk-3}%
seu numerum unciarum ponderis laminae pedis quadrati%
\indexS{uncia}%
\indexS{pondus laminae cupreae}%
\indexS{pes quadratus laminae cupreae}%
\indexS{lamina cuprea}
quadruplicatum%
\lbrack,\rbrack\
esse numerum pedum sphaerae minimae in aere elevabilis.)%
\indexS{sphaera minima elevabilis in aere}%
\indexS{sphaera cuprea}%
\indexS{pes}%
\indexS{aer communis}%
\edlabel{Gerland_1906_007-009_rndbmrk-4}%
\pend%
\vspace{2.0em}%
\pstart%
\noindent%
\centering%
*     *     *
\pend%
\vspace{2.0em}%
\selectlanguage{german}%
\frenchspacing%
\footnotesize%
\pstart%
\noindent%
\textls{Ausführliche Datierungsgründe:}%
\pagelabel{56215_dat}
Die vorliegende undatierte Aufzeichung N.\;\ref{56215} \textit{Aero-nautica}, die lediglich in einem modernen Druck überliefert ist (siehe oben, S.\;\pageref{Gerland_1906_007-009_uebrlf}), setzt die Beschreibung einer imaginären \textit{nave che camini sostentata sopra l'aria} voraus, welche bei F.\,\textsc{Lana\:de\:Terzi}, \textit{Prodromo overo Saggio di alcune inventioni nuove premesso all'arte maestra}, I Teil, cap.\;6 (Brescia 1670, S.\;52--61)\cite{00069} zu lesen und Leibniz bereits im Jahre 1671 bekannt war,
wie dessen Erwähnung in der \textit{Hypothesis physica nova}, §\;23\;\lbrack bis\rbrack\ (\textit{LSB} VI,\;2 N.\;40, S.\;232.30--233.21)\cite{00256} zeigt;
siehe zudem die zur gleichen Zeit entstandenen Auszüge aus Lanas \textit{Prodromo} (\textit{LSB} VIII,\;1 N.\;16).\cite{01103}
In N.\;\ref{56215} beschränkt sich Leibniz darauf, den Durchmesser einer kupfernen Kugel zu bestimmen, welche, von der inneren Luft entleert, in der äußeren Luft gemäß dem hydrostatischem Prinzip schweben sollte.
(Lanas\indexN{\textls{Lana de Terzi}, Francesco 1631\textendash1687} Vorstellung war nämlich, dass eine hinreichende Zahl derartiger Kugeln imstande sei, ein Schiff in die Luft mit zu heben.)
Im ersten Teil der Aufzeichnung (bis \refpassage{Gerland_1906_007-009_scheitel-1}{Gerland_1906_007-009_scheitel-1}) versucht Lebniz, den minimalen Wert des Durchmessers einer solchen Kugel festzulegen, wobei er von den gleichen physikalisch-metrischen Annahmen ausgeht wie Lana selbst (vgl. \textit{Prodromo},\cite{00069} S.\;53; 56):
Ein Kubikfuß Luft wiege anderthalb Unzen;
eine Unze sei der zwölfte Teil eines Pfundes;
ein Kupferblech, das drei Unzen je Quadratfuß wiege, sei für den Bau einer in der Luft schwebenden Kugel sowohl hinreichend dünn wie auch hinreichend fest.
Unter diesen drei Annahmen kommt er auf die Formeln $\displaystyle\frac{66}{7}b^2$ für das Gewicht der kupfernen Kugel und $\displaystyle\frac{11}{14}b^3$ für das Gewicht der eingeschlossenen Luft (mit \textit{b} als Durchmesser), welche eine Verallgemeinerung von Lanas\indexN{\textls{Lana de Terzi}, Francesco 1631\textendash1687} Ergebnissen darstellen.
\pend%
\pstart%
Diese Berechnung wird im zweiten Teil der vorliegenden Aufzeichnung (ab \refpassage{Gerland_1906_007-009_scheitel-2}{Gerland_1906_007-009_scheitel-2}) berichtigt:
Leibniz geht hierbei von der aus seiner Sicht realistischeren Annahme aus, dass eine für den Bau der erwünschten Kugel hinreichend feste Kupferplatte etwa 32 Unzen je Quadratfuß wiegen und somit um etwa das zehnfache dicker sein sollte, als von Lana\indexN{\textls{Lana de Terzi}, Francesco 1631\textendash1687} vorgeschlagen.
Hieraus folgt, dass die Kugel, um sich in der Luft heben zu können, einen Durchmesser von 128 Fuß besitzen sollte, wohingegen Lana\indexN{\textls{Lana de Terzi}, Francesco 1631\textendash1687} zu einem Wert von vierzehn Fuß gekommen war (vgl. \textit{Prodromo},\cite{00069} S.\;56).
Dies zeigt Leibniz zufolge, dass der Gedanke einer luftentleerten kupfernen Kugel, die in der Luft zu schweben vermöge, an seiner praktischen Umsetzung scheitern müsse.
Das allgemeine rechnerische Ergebnis der Untersuchung bleibt jedenfalls -- wie am Ende der Aufzeichnung (\refpassage{Gerland_1906_007-009_schlss-1}{Gerland_1906_007-009_schlss-2}) sowie in einer Randbemerkung im Kopf (\refpassage{Gerland_1906_007-009_rndbmrk-1}{Gerland_1906_007-009_rndbmrk-2}) festgehalten wird --, dass der Durchmesser der schwebenden Kugel mindestens viermal so viele Fuß lang sein müsse, wie Unzen ein Quadratfuß Kupferblech wiege.
(Bemerkenswert ist, dass im zweiten Teil der Aufzeichnung %
Leibniz die vorgefundene Unterteilung des Pfundes in zwölf Unzien aufgibt und stillschweigend durch eine Unterteilung in sechzehn Unzen ersetzt; vgl. \refpassage{Gerland_1906_007-009_crj-3}{Gerland_1906_007-009_crj-4}.)
\pend%
\pstart%
Die Vorlage, an die die vorliegende Aufzeichnung N.\;\ref{56215} \textit{unmittelbar} anknüpft, ist jedoch nicht Lanas \textit{Prodromo}\cite{00069}, sondern offenbar die von Philipp Lohmaier\indexN{\textls{Lohmeier} (Lohmeierus, Lomeierus), Philipp 1648\textendash1680} betreute Dissertation \textit{Exercitatio physica de artificio navigandi per aerem} (Rinteln 1676),\cite{01422} in welcher Lanas Gedanke eines Luftschiffes stillschweigend übernommen und erörtert wird.
Weder Lana\indexN{\textls{Lana de Terzi}, Francesco 1631\textendash1687} noch Lohmeier werden in N.\;\ref{56215} ausdrücklich erwähnt.
Die Sätze, mit denen am Anfang der Aufzeichnung die physikalischen Grundannahmen des Gedankensexperiments erklärt werden, geben jedoch nahezu wörtlich die entsprechenden Textpassagen der Dissertation wieder:
Vgl. N.\;\ref{56215}, \refpassage{Gerland_1906_007-009_hestz-1}{Gerland_1906_007-009_hestz-2} mit \textit{Exercitatio physica},\cite{01422} §\;8, S.\;7f. (\textit{Deprehendunt ... pondus illud} \lbrack\textit{= aeris}\rbrack\ \textit{accurati observatores tale, ut unus pes cubicus aeris aequet unciam unam cum dimidia ...})
sowie N.\;\ref{56215}, \refpassage{Gerland_1906_007-009_xtd-1}{Gerland_1906_007-009_crj-2} mit \textit{Exercitatio physica},\cite{01422} §\;21, S.\;14 (\textit{Denique suppono ... parari posse e cupro laminas aliquas satis adhuc consistentes globisque inde conficien}\lbrack\textit{d}\rbrack\textit{is aptas, quarum longitudo et latitudo unius pedis sit, et tamen pondere non excedant tres uncias ...})
und mit §\;30, S.\;18 (\textit{... assumo ... parari posse e cupro laminas aliquas pedem unum longas et latas, satis adhuc consistentes et pondere suo tres uncias non excedentes ...}).
Ferner lässt sich ein unmittelbarer Zusammenhang zwischen der Aufzeichnung N.\;\ref{56215} und Lohmeiers\indexN{\textls{Lohmeier} (Lohmeierus, Lomeierus), Philipp 1648\textendash1680} Brief an Leibniz vom 23.\;August (2.\;September) 1679 ermitteln, in welchem die \textit{Exercitatio physica}\cite{01422} ausführlich erwähnt und deren Verknüpfung mit Lanas\indexN{\textls{Lana de Terzi}, Francesco 1631\textendash1687} Gedanken eines schwebenden Schiffes berührt werden (\textit{LSB} III,\;2 N.\;336, S.\;815.7--13).\cite{01425}
Dort bittet Lohmeier\indexN{\textls{Lohmeier} (Lohmeierus, Lomeierus), Philipp 1648\textendash1680} seinen Korrespondenten darum, mögliche Beobachtungen über das Thema mitteilen zu wollen, welche in die geplante zweite Auflage der \textit{Exercitatio physica}\cite{01422} aufgenommen werden könnten (\textit{LSB} III,\;2 N.\;336, S.\;815.15--816.2; die zweite Auflage der Dissertation erschien 1679 in Wittenberg, jedoch ohne inhaltliche Änderungen gegenüber der ersten Auflage).
In einer am Briefrand verfassten Anmerkung zu dieser Stelle hat Leibniz mit nahezu gleichem Wortlaut das rechnerische Ergebnis der in N.\;\ref{56215} durchgeführten Untersuchung (\refpassage{Gerland_1906_007-009_schlss-1}{Gerland_1906_007-009_schlss-2}; \refpassageohne{Gerland_1906_007-009_rndbmrk-3}{Gerland_1906_007-009_rndbmrk-4}) festgehalten: \textit{Numerus unciarum laminae pedis quadrati quadruplicatus dat numerum pedum diametri sphaerae} (\textit{LSB} III,\;2 N.\;336, S.\;815.20\,f.).
Dieser Zusammenhang ermöglicht die Datierung der vorliegenden Aufzeichnung.
\pend%
\pstart%
Leibniz beantwortete Lohmeiers Schreiben mit zwei Briefen, heute beide verschollen: Der erste (\textit{LSB} III,\;2 N.\;354)\cite{01426} wurde wahrscheinlich zwischen Ende Oktober und Anfang November 1679, der zweite (\textit{LSB} III,\;2 N.\;369)\cite{01427} im folgenden Dezember versendet.
Lohmeiers\indexN{\textls{Lohmeier} (Lohmeierus, Lomeierus), Philipp 1648\textendash1680} Erwiderung vom 24.\;Dezember 1679 (3.\;Januar 1680) lässt sich entnehmen, dass Leibniz wohl in beiden seiner Antworten auf Lohmeiers Bitte eingegangen war und seine \textit{mens de artificio aeronautico} mitgeteilt hatte (\textit{LSB} III,\;3 N.\;1, S.\;43.11--13).\cite{01428}
Die Vermutung liegt somit nahe, dass Leibniz, um Lohmeiers Brief vom 23.\;August (2.\;September) zu beantworten, die \textit{Exercitatio physica}\cite{01422} gelesen und im Anschluss die Aufzeichnung N.\;\ref{56215} verfasst hat.
In seinen Briefen an Lohmeier hat er wahrscheinlich auf die Ergebnisse seiner Berechnung in N.\;\ref{56215} Bezug genommen.
Unter dieser Annahme ist die vorliegende Aufzeichnung auf den Zeitraum zwischen September und Dezember 1679 zu datieren.
\pend%
\pstart%
Es ist jedoch möglich, dass Leibniz bereits vor Lohmeiers Brief vom 23.\;August (2.\;September) 1679 die \textit{Exercitatio physica}\cite{01422} gelesen und die Aufzeichnung N.\;\ref{56215} verfasst hatte.
Hierfür spricht, dass er bereits vor Mitte August 1679 eine weitere, von Lohmeier betreute Dissertation, die \textit{Exercitationum physicarum de paradoxis gravitatis et levitatis prima} (Rinteln 1678),\cite{01429} gründlich gelesen hatte, wie dies sich dem letztgenannten Brief entnehmen lässt (\textit{LSB} III\;2, N.\;336, S.\;817.6--15, wo auf einen Zweifel eingegangen wird, den Leibniz in seinem -- verschollenen -- Schreiben an Lohmeier\indexN{\textls{Lohmeier} (Lohmeierus, Lomeierus), Philipp 1648\textendash1680} von Mitte August 1679 %
geäußert haben muss).
Daher ist nicht unwahrscheinlich, dass er zur selben Zeit oder gar früher auch Zugang zu einem Exemplar der \textit{Exercitatio physica}\cite{01422} gehabt und die in N.\;\ref{56215} überlieferten Berechnungen angestellt hatte, auf die er später, nach dem Empfang von Lohmeiers Brief vom 23.\;August (2.\;September) 1679,\cite{01425} zurückgriff, um diesen zu beantworten.
Unter dieser Annahme kann die vorliegende Aufzeichnung frühestens Ende 1676 verfasst worden sein.
Denn erst nach seiner Rückkehr aus Paris und nach seiner Ankunft in Hannover Mitte Dezember 1676 konnte Leibniz von Lohmeiers\indexN{\textls{Lohmeier} (Lohmeierus, Lomeierus), Philipp 1648\textendash1680} \textit{Exercitatio physica},\cite{01422} die nicht vor März desselben Jahres gedruckt worden war, Kenntnis nehmen.
Eine frühere, auf Leibnizens erste Lektüre von Lanas \textit{Prodromo}\cite{00069} zurückgehende Entstehung von N.\;\ref{56215} ist durch die oben erwähnte textuelle Abhängigkeit der Aufzeichnung von der \textit{Exercitatio physica}\cite{01422} ausgeschlossen.
Hieraus ergibt sich als Gesamtdatierung der vorliegenden Aufzeichnung N.\;\ref{56215} die Zeitspanne zwischen der zweiten Hälfte Dezember 1676 und Dezember 1679.
\pend%
\selectlanguage{latin}%
\frenchspacing%
\normalsize%
\newpage%
